\section*{Discussion}

We developed two fine-mapping methods---hierarchical belief
propagation (HBP) and graph-augmented
fine-mapping (GAFM)---that represent multi-omics biology as a typed factor
graph rather than as a flat per-variant annotation vector. HBP
matches state-of-the-art Bayesian accuracy at 5--40$\times$ the
speed under strong signal; GAFM wins decisively against SuSiE at
weak signal with tissue-specific eQTL priors (27--2 head-to-head,
$p < 10^{-5}$). Both methods are PIP-calibrated with 0\% null
false-positive rate, and both scale from 1000 Genomes to Pan-UK
Biobank across four ancestries without algorithmic modification.

\textbf{Graph-native fine-mapping is structurally distinct from
flat-prior fine-mapping, not merely faster.} The head-to-head
benchmarks quantify this distinction on the
\emph{representational} axis. When the prior information supplied
to SuSiE matches what HBP/GAFM use (Polyfun-proxy with eQTL-derived
\texttt{prior\_weights}), SuSiE and GAFM converge on rank-\#1
accuracy; GAFM wins only on the $26 \times$ speed margin.  But when
the graph captures relational structure that a per-variant weight
vector cannot---tissue-specific eQTL enrichment, PPI coupling on
the gene layer, pathway co-occurrence---the graph prior wins not
only on speed but on rank-\#1 rate (27--2), mean rank (1.65 vs
2.84) and calibrated PIP. The implication is categorical: as
multi-omics resources continue to densify with tissue-resolved
and context-specific annotations, the advantage of relational
priors will widen, not narrow. This is the regime into which the
next decade of annotation coverage is entering. The
intervention experiment in
\S``\nameref{sec:heterogeneity}'' sharpens this further: density alone is
insufficient. The dense uniform GTEx+STRING human cache yielded a
null intervention (HBP exactly equal to GAFM on $321/321$ loci), but
adding a coarse heterogeneous regulatory-element layer
(ENCODE cCRE) flipped the HBP-tighter rate from $0\%$ to $88\%$
on the same loci. The right axis to optimise is per-variant
discriminative information, not annotation count.

\textbf{The representational shift is coupled to a computational
shift.} The second axis of distinction---less visible in the
numbers but central to the method's viability---is the data model.
Matrix- and file-based pipelines treat annotations as
per-variant scalars because that is what the data model cheaply
supports: reconstructing the full multi-omics neighbourhood of a
variant on demand from BED, GTF, VCF and TSV files is
prohibitively slow at biobank scale. A graph
database\cite{robinson2015graph} moves this neighbourhood into
primary storage: typed nodes and labelled edges allow the
variant~$\rightarrow$~gene~$\rightarrow$~tissue~$\rightarrow$~pathway traversal
that HBP needs to execute thousands of times per locus to return
in milliseconds. Without this computational substrate, HBP's
three-layer message passing would be formally expressible but
operationally unusable; with it, the factor graph can be queried,
sliced and updated as cheaply as SuSiE queries its correlation
matrix. The shift we report therefore has two coupled
halves: the representational move from flat annotation vectors to
relational priors, and the computational move from
file-and-matrix pipelines to graph-database substrates that make
those relational priors tractable. Either half alone would be
insufficient.

\textbf{When SuSiE and FINEMAP remain the right tool.} Under
pure statistical conditions---strong signal, random causal
variant, no informative annotations---SuSiE\cite{wang2020simple,zou2022fine}
and FINEMAP\cite{benner2016finemap} retain a small advantage in
mean rank\cite{hutchinson2020improving} and should remain the
default.
GraphGWAS is a complement, not a replacement: it targets the
weak-signal and annotation-rich regime that matters most for
post-GWAS discovery, and it targets workloads where per-locus
speed dominates wall-clock time (many loci, iterative
exploration, interactive analysis). Similarly, SBayesRC and
related polygenic methods address a different problem
(genome-wide credible-set construction at $N \gtrsim 10^5$);
HBP/GAFM and SBayesRC are complementary within a workflow, with
SBayesRC surfacing leads and HBP/GAFM refining them per-locus.

\textbf{Non-model species --- crops, livestock, and model
organisms --- removing the annotation-weight bottleneck.}\label{disc:nonmodel} Flat-prior fine-mappers depend
on a reliable per-variant annotation weight vector
$w \in \mathbb{R}^n$, typically produced by stratified LD-score
regression over curated functional categories
(Baseline 2.2 in humans uses 96
categories)\cite{finucane2015partitioning,buliksullivan2015ld}.
This pipeline has three hard prerequisites: (i) a large well-
powered GWAS ($N \gtrsim 20{,}000$) to estimate stable
per-category heritability enrichment; (ii) pre-curated and
biologically meaningful functional categories; and (iii) an LD
reference matching the target cohort's population. All three
are met in human biobank analyses, and none of them is reliably
met in GWAS on crops, livestock, and other non-model species. The 3{,}024-
accession 3K~Rice~Genomes panel\cite{wang2018genomic}, the
canonical crop-genomics resource for fine-mapping validation,
has $N$ an order of magnitude too small for stable S-LDSC
enrichment estimation, no species-specific Baseline-equivalent
annotation catalogue, and 12--15 cross-subspecies LD strata
(GJ versus XI $F_{ST} \approx 0.5$) that violate the single-
ancestry assumption of LD-score regression. Comparable gaps
exist for maize, wheat, soybean and essentially every non-human
species of agricultural or ecological relevance: the scalar
annotation weight vector cannot be constructed without heavy
manual curation that is itself unvalidated. Graph-native
fine-mapping sidesteps this bottleneck entirely because it never
demands a harmonised weight vector. Annotation layers register
as typed edges when data exists: for rice, this includes the
RAP-DB / MSU gene models\cite{frankish2023gencode},
Gramene / Reactome Plant
pathways\cite{gillespie2022reactome}, and two complementary
rice-specific protein--protein interaction
resources---PRIN\cite{gu2011prin} (experimental and interolog-
predicted) and RicePPINet\cite{liu2017riceppinet} (708{,}819
interactions across 16{,}895 proteins, ML-predicted from 11
structural, phylogenetic and co-expression features). The
rice-specific PPI networks are nearly two orders of magnitude
denser than the STRING cross-species subset for rice, and are
the kind of species-tailored resource that emerges once a
community invests in a model organism. Tissue-expression
atlases (RiceXPro) register as proxy-eQTL edges where cis-eQTLs
themselves are absent. Missing edges are simply missing, not
re-imputed as zero weights. The graph prior degrades
gracefully with annotation depth---HBP's factor graph reduces to
variant~$\rightarrow$~gene propagation when only gene models
exist, and each additional edge type (pathway, PPI, tissue-
specific eQTL) contributes additively. For crop and livestock breeders and other non-model species
researchers whose GWAS panels hit the
statistical-power and annotation-curation walls at which
flat-prior methods require costly domain-specific pipelines,
this is a categorical removal of a recurring bottleneck rather
than a marginal refinement. The 3{,}000 Rice Genomes scan
reported here recovers 24 of 72 fine-mapped leads (33\%) within
250~kb of a Ren, Ding \& Qian (2023)\cite{ren2023rice}-curated
causal gene including \emph{Rc} (4.5~kb), \emph{TGW6} (12.3~kb),
\emph{GAD1/RAE2}, \emph{An-1} and \emph{BG1} --- a concrete
test of the portability claim on the canonical crop-genomics
benchmark species. Although the present paper benchmarks the
cross-species claim only on yeast, \emph{Arabidopsis}, rice, and
human data, the same annotation-curation bottleneck dominates
livestock GWAS --- cattle, pig, chicken and sheep panels
increasingly assembled by FarmGTEx-style consortia --- and we
expect the graph-native pipeline to transfer there with similar
gains as species-specific PPI and tissue-eQTL resources mature.

\textbf{Limitations.} Several limitations bound the present
claims, with three pre-registered placebo / null-permutation /
sensitivity experiments now reported in Sup Note~S2. (i)~A
placebo-prior control on a chr~22 strong-signal simulation
(30 replicates, $\beta{=}0.5$, $h^2{=}0.10$) shows that the gain
from adding a non-uniform prior (10\,\%\,$\to$\,13\,\%
HBP-tighter rate over baseline) is \emph{fully captured} by a
permuted prior with the same marginal distribution: informative
cCRE prior 13\,\%, placebo permuted prior 13\,\% (5 permutations,
zero variance). In this simulation regime, with randomly-placed
causal variants, biological structure adds nothing on top of
heterogeneity per se. The original 0\,\%\,$\to$\,88\,\% flip on
321 Pan-UK Biobank lead loci (Table~1) therefore depends on
enrichment of real disease-associated leads in cCRE-class regions
(a known phenomenon) rather than on the graph prior providing
variant-specific biological information beyond what any
non-uniform prior would produce. We retain the Table~1
intervention result but reinterpret it in this light: the typed
factor graph is the necessary substrate for incorporating
non-uniform priors into HBP/GAFM; the empirical superiority of
informative-cCRE over a permuted-marginal placebo on real
biobank leads remains to be demonstrated. (ii)~The 250\,kb-window cross-species
validation rate is reported without a species-matched
random-position null; without that null, the per-species rates
quoted in §``\nameref{sec:multispecies}'' should be read as
\emph{lower bounds} on overlap, not as fine-mapping accuracy. (iii)~The
unified six-method head-to-head benchmark in Figure~\ref{fig:baselines}
uses 30 replicates per scenario; pairwise rank-\#1 differences of
three percentage points are within Wilson 95\% CIs and should be
read as ``no significant difference at $\alpha=0.05$'' rather than
as equivalence. (iv)~HBP and GAFM are reported with default
hyperparameters ($\alpha$, $\lambda$, $T$ for HBP; $\alpha$,
$\tau_L$ for GAFM); a one-at-a-time sensitivity analysis is
forthcoming. (v)~Pan-UKB sample size is heavily dominant in EUR
($N=420{,}531$ vs $\leq 8{,}876$ in the three other ancestries),
so cross-ancestry rank-\#1 comparisons confound ancestry effects
with statistical power.

Beyond these immediate gaps, pathway annotations remain sparse on
chromosome~22 in our current graph (5 pathways with variants in
the weak-signal test windows); the full benefit of pathway-layer
propagation will be realised with denser Reactome\cite{gillespie2022reactome}
and KEGG\cite{kanehisa2023kegg} coverage. Belief propagation on loopy factor graphs is
approximate; HBP's softmax caps posterior mass at $\approx$0.7
per variant by design --- HBP and GAFM should therefore be read on
different posterior scales (HBP for credible-set narrowness; GAFM
for absolute calibration up to PIP\,=\,1). Our Pan-UKB demonstration uses
1KG-matched LD as a stand-in while the \texttt{hadoop-aws} Spark
configuration needed to read Pan-UKB in-sample LD BlockMatrices
directly is completed---a software-plumbing task, not a
methodological one, that we expect will tighten credible sets
for small-$N$ ancestries further. Finally, we have not yet
applied GraphGWAS at full UK Biobank scale (direct analysis of
individual-level 500~K genotypes); the hybrid BGEN~+~graph
architecture is proven on 1000 Genomes and on Pan-UKB sumstats,
but full-UKB application requires dataset access.

\textbf{The graph substrate generalises beyond fine-mapping.}
The methodological move---carrying relational biological
structure through inference rather than collapsing it to a
scalar---applies to problems adjacent to
fine-mapping\cite{ulirsch2021fine,schaid2018genome}. Cross-trait
pleiotropy becomes a graph centrality statistic when multiple
GWAS studies write their \texttt{AssociationResult} nodes into a
shared graph; Mendelian randomisation becomes traversal of a
typed causal DAG; label-propagation diffusion of phenotype
labels from sparsely-phenotyped cohorts becomes a native graph
operation. Higher-order epistasis---pair and triplet
interactions---is particularly naturally expressed as
subgraph-enumeration on the LD-pruned graph
(Supplementary Fig.~S1), with a search-space reduction of
42{,}000$\times$ on chromosome~22 common variants (Theorem~1).
Learned message passing (graph neural networks on the same
factor graph) is the natural lift of HBP and will follow this
work. Association-rule mining over the same typed factor graph ---
combinations of variants, tissues, pathways and regulatory contexts
that co-occur in cases more than chance --- is a further natural
follow-up, since the graph substrate exposes multi-relational
patterns directly to rule-mining algorithms; the LD-pruned
independent-set construction we develop for LPCE generalises
straightforwardly to the higher-order rules this would produce.
The per-locus fine-mapping result is the proof of
concept; the broader claim---that post-GWAS analysis, conducted
on a data model that matches the biology, is more naturally a
graph computation than a matrix computation---is the reframing
this work establishes.
